% =====================================================================
%  DataPilot — Rapport de Projet de Fin d'Études
%  Compilation : pdflatex (x2 pour la table des matières / références)
%      pdflatex DataPilot_PFE_Rapport.tex
%      pdflatex DataPilot_PFE_Rapport.tex
% =====================================================================
\documentclass[11pt,a4paper]{report}

\usepackage[utf8]{inputenc}
\usepackage[T1]{fontenc}
\usepackage{lmodern}
\usepackage[french]{babel}
\usepackage[a4paper,margin=2.5cm]{geometry}
\usepackage{graphicx}
\usepackage{xcolor}
\usepackage{booktabs}
\usepackage{array}
\usepackage{enumitem}
\usepackage{fancyhdr}
\usepackage{tikz}
\usetikzlibrary{shapes.geometric, arrows.meta, positioning}
\usepackage{listings}
\usepackage[hidelinks]{hyperref}

% --- Couleurs de marque EY ---
\definecolor{eyyellow}{HTML}{FFE600}
\definecolor{eycharcoal}{HTML}{2E2E38}
\definecolor{eypurple}{HTML}{3D108A}
\definecolor{lightgray}{HTML}{F2F2F2}

% --- En-têtes / pieds de page ---
\pagestyle{fancy}
\fancyhf{}
\setlength{\headheight}{14pt}
\fancyhead[L]{\small\textsc{DataPilot}}
\fancyhead[R]{\small Évaluation de la maturité data bancaire}
\fancyfoot[C]{\thepage}
\renewcommand{\headrulewidth}{0.4pt}

% --- Style des listings (code / SQL) ---
\lstdefinestyle{code}{
  basicstyle=\ttfamily\footnotesize,
  keywordstyle=\color{eypurple}\bfseries,
  commentstyle=\color{gray}\itshape,
  stringstyle=\color{teal},
  backgroundcolor=\color{lightgray},
  frame=single, rulecolor=\color{gray!40},
  breaklines=true, showstringspaces=false,
  numbers=left, numberstyle=\tiny\color{gray}, numbersep=8pt,
  captionpos=b, tabsize=2,
}
\lstset{style=code}

\renewcommand{\arraystretch}{1.3}

% =====================================================================
\begin{document}

% ----------------------------- PAGE DE GARDE -----------------------------
\begin{titlepage}
\centering
\vspace*{1cm}

{\Huge\bfseries\color{eycharcoal} EY \;}\rule{2pt}{1.5em}\; {\Large DataPilot}\\[0.3cm]
{\large\itshape Data Maturity Steering Tool}\\[2cm]

\rule{\linewidth}{0.5mm}\\[0.5cm]
{\LARGE\bfseries Conception et réalisation d'un outil web d'évaluation\\[0.2cm]
de la maturité de la gestion des données bancaires}\\[0.4cm]
\rule{\linewidth}{0.5mm}\\[1.5cm]

{\large
\begin{tabular}{rl}
\textbf{Projet de Fin d'Études} & 2026\\[0.3cm]
\textbf{Organisme d'accueil} & EY Advisory Tunisie\\[0.3cm]
\textbf{Domaine} & Gouvernance \& maturité des données\\[0.3cm]
\textbf{Conformité} & BCT Circulaire N\textsuperscript{o}2025-08, BCBS 239\\
\end{tabular}}

\vfill
{\large Élaboré par : \textbf{Aziz Dardouri}}\\[0.3cm]
{\normalsize Encadré par : votre encadrant académique \& votre encadrant professionnel}\\[1cm]
{\small Année universitaire 2025--2026}

\end{titlepage}

% ----------------------------- RÉSUMÉ -----------------------------
\chapter*{Résumé}
\addcontentsline{toc}{chapter}{Résumé}

\textbf{DataPilot} est une application web qui permet aux banques tunisiennes
d'évaluer objectivement la maturité de leur gestion des données. À partir d'un
questionnaire structuré de 47 indicateurs répartis sur 5 dimensions pondérées,
l'outil calcule un score de maturité aligné sur les référentiels CMMI et
Gartner, le confronte aux exigences réglementaires de la Banque Centrale de
Tunisie (Circulaire N\textsuperscript{o}2025-08) et au standard BCBS~239, puis
génère une analyse d'écarts ainsi qu'une feuille de route d'amélioration
priorisée. La solution intègre une authentification sécurisée sur invitation,
une gestion des rôles, un back-office d'administration permettant d'éditer
intégralement le questionnaire, et l'export de rapports PDF de qualité
professionnelle.\\

\textbf{Mots-clés :} maturité des données, gouvernance des données, conformité
bancaire, BCT, BCBS 239, CMMI, React, Supabase.\\

\noindent\textbf{Abstract.} \textit{DataPilot is a web application enabling
Tunisian banks to objectively assess their data-management maturity. From a
structured questionnaire of 47 indicators across 5 weighted dimensions, it
computes a CMMI/Gartner-aligned maturity score, benchmarks it against the
Central Bank of Tunisia regulation (Circular 2025-08) and the BCBS~239 standard,
and produces a gap analysis and a prioritised improvement roadmap. The solution
features invite-only authentication, role management, a full content-editing
back-office, and professional PDF report export.}

\tableofcontents

% ===================================================================
\chapter{Introduction générale}

\section{Contexte}
La donnée est devenue un actif stratégique du secteur bancaire. Sa bonne
gouvernance conditionne la qualité du reporting réglementaire, la maîtrise des
risques et la capacité d'innovation. En Tunisie, la \emph{Circulaire BCT
N\textsuperscript{o}2025-08} impose aux établissements de crédit un ensemble
d'obligations en matière d'organisation, de qualité et de traçabilité des
données, en cohérence avec le standard international \emph{BCBS~239} relatif à
l'agrégation des données de risque.

\section{Problématique}
Malgré ces exigences, les banques manquent souvent d'un outil structuré pour
\textbf{mesurer objectivement} leur niveau de maturité, \textbf{identifier} leurs
écarts par rapport à la réglementation et \textbf{prioriser} leurs actions de
remédiation. Les évaluations restent fréquemment informelles, non reproductibles
et difficilement comparables dans le temps.

\section{Objectifs}
Le présent projet vise à concevoir et réaliser \textbf{DataPilot}, un outil web
qui :
\begin{itemize}[noitemsep]
  \item structure l'évaluation autour d'un référentiel d'indicateurs pondérés ;
  \item calcule un score de maturité reproductible et traçable ;
  \item mesure la conformité réglementaire (BCT, BCBS~239) ;
  \item produit une analyse d'écarts et une feuille de route d'amélioration ;
  \item génère des rapports PDF exportables et de qualité professionnelle ;
  \item offre un back-office d'administration sécurisé et paramétrable.
\end{itemize}

\section{Organisation du rapport}
Le chapitre~2 présente le cadre du projet et l'état de l'art. Le chapitre~3
détaille la méthodologie d'évaluation. Le chapitre~4 décrit la conception et
l'architecture. Le chapitre~5 expose la réalisation. Le chapitre~6 présente les
fonctionnalités obtenues. Le chapitre~7 conclut et ouvre des perspectives.

% ===================================================================
\chapter{Cadre du projet et état de l'art}

\section{Organisme d'accueil}
Le projet est réalisé au sein d'\textbf{EY Advisory Tunisie}, entité de conseil
du cabinet EY, intervenant notamment sur la transformation data et la conformité
réglementaire du secteur financier.

\section{La maturité des données}
La \emph{maturité des données} désigne le degré auquel une organisation gère ses
données de manière structurée, fiable et orientée valeur. Elle est usuellement
mesurée sur une échelle de cinq niveaux, ici alignée sur deux référentiels
reconnus :
\begin{itemize}[noitemsep]
  \item \textbf{CMMI} : Initial, Emerging, Defined, Managed, Optimized ;
  \item \textbf{Gartner} : Unaware, Aware, Active, Effective, Transformative.
\end{itemize}

\section{Cadre réglementaire}
\begin{itemize}[noitemsep]
  \item \textbf{Circulaire BCT N\textsuperscript{o}2025-08} : obligations
  d'organisation, de qualité, de propriété et de traçabilité des données.
  \item \textbf{BCBS~239} : principes d'agrégation et de reporting des données de
  risque (notamment le Principe~2 — traçabilité / lignage des données).
\end{itemize}

\section{Positionnement}
Contrairement aux audits ponctuels, DataPilot fournit un instrument
\textbf{reproductible} et \textbf{auto-administrable}, dont le référentiel peut
évoluer dans le temps sans dépendre d'un prestataire externe.

% ===================================================================
\chapter{Méthodologie d'évaluation}

\section{Les cinq dimensions}
L'évaluation s'articule autour de cinq dimensions pondérées, déclinées en
sous-dimensions et en indicateurs.

\begin{table}[h]
\centering
\caption{Dimensions, pondérations et nombre d'indicateurs (configuration par défaut).}
\begin{tabular}{clcc}
\toprule
\textbf{Dim.} & \textbf{Intitulé} & \textbf{Poids} & \textbf{Indicateurs}\\
\midrule
D1 & Gouvernance & 25\,\% & 12 \\
D2 & Qualité des données & 20\,\% & 11 \\
D3 & Architecture \& Accès & 20\,\% & 8 \\
D4 & Analytique \& Outils & 20\,\% & 8 \\
D5 & Compétences \& Culture (proxy) & 15\,\% & 8 \\
\midrule
\multicolumn{2}{r}{\textbf{Total}} & 100\,\% & \textbf{47}\\
\bottomrule
\end{tabular}
\end{table}

\section{Grille de notation}
Chaque indicateur est noté de 1 à 5 selon une grille (\emph{rubric}) décrivant
précisément l'état attendu à chaque niveau. Le score d'une sous-dimension est la
moyenne de ses indicateurs ; le score d'une dimension la moyenne de ses
sous-dimensions ; et le \textbf{score global} la moyenne \emph{pondérée}
(normalisée) des dimensions effectivement renseignées :
\[
S_{global} = \frac{\sum_{d} w_d \cdot S_d}{\sum_{d} w_d}, \qquad
S_d \in [1,5].
\]

\section{Règles métier}
Afin de garantir la crédibilité du diagnostic, plusieurs règles sont appliquées
(tableau~\ref{tab:rules}).

\begin{table}[h]
\centering
\caption{Règles de notation appliquées par le moteur d'évaluation.}
\label{tab:rules}
\begin{tabular}{p{4.5cm} p{9.5cm}}
\toprule
\textbf{Règle} & \textbf{Description}\\
\midrule
Plafonnement par preuve & Un score $\geq 3$ non justifié par une preuve
documentée est plafonné à 2/5.\\
Limite de saut & Au plus 20\,\% des indicateurs d'une dimension peuvent être
ignorés.\\
Indicateurs BCT & Les indicateurs réglementaires (BCT) sont obligatoires et ne
peuvent pas être ignorés.\\
Dimension proxy & La dimension \og Compétences \& Culture \fg{} est évaluée par
des indicateurs \emph{proxy} (signaux objectifs) plutôt que par auto-déclaration.\\
\bottomrule
\end{tabular}
\end{table}

\section{Conformité réglementaire}
Un sous-ensemble d'indicateurs est marqué \textbf{BCT}. Un indicateur est jugé
\emph{conforme} si son score effectif est $\geq 3$. Le taux de conformité
détermine l'exposition réglementaire : \emph{Faible} ($\geq 80\,\%$),
\emph{Moyenne} (50--79\,\%) ou \emph{Élevée} ($< 50\,\%$).

% ===================================================================
\chapter{Conception et architecture}

\section{Vue d'ensemble}
DataPilot adopte une architecture web moderne : une application monopage (SPA)
React communiquant avec une base de données managée Supabase (PostgreSQL +
authentification), complétée de fonctions serverless pour les opérations
sensibles (génération IA de la feuille de route, invitation d'utilisateurs).

\begin{figure}[h]
\centering
\begin{tikzpicture}[
  node distance=1.1cm and 1.8cm,
  box/.style={rectangle, rounded corners, draw=eycharcoal, very thick,
    fill=lightgray, text width=3.2cm, align=center, minimum height=1.3cm},
  db/.style={cylinder, shape border rotate=90, aspect=0.25, draw=eypurple,
    very thick, fill=eyyellow!30, text width=2.6cm, align=center, minimum height=1.4cm},
  arr/.style={-{Latex[length=2.5mm]}, thick}
]
\node[box] (browser) {Navigateur\\\textbf{React SPA}\\(Vite, Tailwind, Zustand)};
\node[box, right=of browser] (api) {Fonctions\\\textbf{Serverless}\\/api/invite, /api/roadmap};
\node[db, right=of api] (supa) {\textbf{Supabase}\\PostgreSQL\\+ Auth + RLS};
\node[box, below=of api] (ai) {Service IA\\(endpoint compatible\\OpenAI)};

\draw[arr] (browser) -- node[above,font=\scriptsize]{HTTPS} (api);
\draw[arr] (browser.south) to[bend right=15] node[below,font=\scriptsize]{anon key} (supa.south west);
\draw[arr] (api) -- node[above,font=\scriptsize]{service\_role} (supa);
\draw[arr] (api) -- (ai);
\end{tikzpicture}
\caption{Architecture technique de DataPilot.}
\end{figure}

\section{Pile technologique}
\begin{itemize}[noitemsep]
  \item \textbf{Front-end} : React~19, Vite, Tailwind CSS, Zustand (gestion
  d'état), React Router, Recharts (visualisations), \texttt{react-to-print}.
  \item \textbf{Back-end} : Supabase (PostgreSQL managé, région UE,
  authentification), fonctions serverless (style Vercel) avec équivalent en
  middleware de développement Vite.
  \item \textbf{IA} : génération de recommandations via un endpoint compatible
  OpenAI (Gemini, palier gratuit), avec repli sur des recommandations statiques.
\end{itemize}

\section{Justification du choix de Supabase}
Supabase regroupe en un seul service la base de données, l'authentification et la
sécurité au niveau ligne (RLS). Gratuit, open-source et donc \textbf{portable}
(un déploiement on-premise reste une simple reconfiguration), il convient à un
projet de fin d'études tout en demeurant industrialisable.

\section{Modèle de données}
Trois tables principales structurent le domaine (en plus des tables
d'authentification gérées par Supabase) :
\begin{itemize}[noitemsep]
  \item \texttt{profiles}(id, email, full\_name, \textbf{role}, created\_at) ;
  \item \texttt{dimensions}(code, name, weight, color, proxy, description,
  sort\_order) ;
  \item \texttt{indicators}(id, dim, sub, sub\_name, bct, q, hint,
  \textbf{rubric} (jsonb), sort\_order).
\end{itemize}
Les sous-dimensions sont dérivées des indicateurs, ce qui permet une gestion
souple du contenu (ajout/suppression de dimensions, sous-dimensions et
indicateurs).

\section{Sécurité}
La sécurité repose sur plusieurs piliers :
\begin{itemize}[noitemsep]
  \item \textbf{Accès sur invitation} : l'inscription publique est désactivée ;
  \item \textbf{Rôles} : \texttt{admin} et \texttt{analyst} ;
  \item \textbf{Row Level Security (RLS)} : chaque accès passe par une politique
  explicite ; la lecture du questionnaire est ouverte aux utilisateurs
  authentifiés, l'écriture est réservée aux administrateurs ;
  \item \textbf{Clé \texttt{service\_role}} : strictement cantonnée au serveur,
  jamais exposée au navigateur.
\end{itemize}

La fonction \texttt{is\_admin()} (en \emph{security definer}) évite la récursion
des politiques RLS sur la table \texttt{profiles} :

\begin{lstlisting}[language=SQL, caption={Fonction de vérification du rôle administrateur.}]
create or replace function public.is_admin()
returns boolean language sql stable
security definer set search_path = public as $$
  select exists (
    select 1 from public.profiles
    where id = auth.uid() and role = 'admin'
  );
$$;
\end{lstlisting}

% ===================================================================
\chapter{Réalisation}

Le développement a été conduit de manière incrémentale, en phases livrables.

\section{Phase 1 — Authentification et rôles}
Mise en place de l'authentification Supabase (e-mail / mot de passe), du modèle
de rôles \texttt{admin}/\texttt{analyst}, de la protection des routes
(\emph{guards} \texttt{RequireAuth} et \texttt{RequireAdmin}) et d'un mécanisme
de dégradation gracieuse lorsque le backend n'est pas configuré.

\section{Phase 2 — Back-office d'administration}
\begin{itemize}[noitemsep]
  \item \textbf{Questionnaire éditable} : édition des questions, des guides, du
  drapeau BCT, des grilles de notation (1--5), des pondérations et des
  descriptions de dimensions ; \textbf{ajout/suppression} d'indicateurs, de
  sous-dimensions et de dimensions.
  \item \textbf{Gestion des utilisateurs} : liste des comptes, changement de
  rôle, et \textbf{invitation par e-mail} directement depuis l'application
  (via une fonction serverless utilisant la clé \texttt{service\_role} après
  vérification du rôle administrateur de l'appelant).
\end{itemize}

Le contenu est chargé depuis Supabase au démarrage, avec repli automatique sur le
questionnaire intégré si la base est vide ou inaccessible : l'application reste
ainsi toujours fonctionnelle.

\section{Génération des rapports PDF}
Trois rapports brandés EY sont produits via le moteur d'impression natif du
navigateur (rendu vectoriel net) :
\begin{itemize}[noitemsep]
  \item \textbf{Rapport de maturité} (Results) : résumé exécutif, radar des
  dimensions, scores pondérés, grille d'interprétation, statut de conformité
  BCT, \textbf{feuille de route} et annexe méthodologique ;
  \item \textbf{Analyse d'écarts} (Gap Analysis) : matrice effort/impact et
  feuille de route en trois phases, avec actions assistées par IA ;
  \item \textbf{Rapport de conformité} (Compliance) : synthèse des actions
  requises, indicateurs non conformes / en attente avec piste de remédiation,
  et tableau détaillé des indicateurs BCT.
\end{itemize}

\section{Feuille de route assistée par IA}
Pour chaque sous-dimension en écart, une fonction serverless interroge un modèle
de langage afin de produire des recommandations concrètes et priorisées, en
traitant en priorité les obligations réglementaires (BCT). En l'absence de
configuration IA, des recommandations standard prennent le relais.

% ===================================================================
\chapter{Fonctionnalités et résultats}

\section{Parcours utilisateur}
Le parcours se déroule en cinq étapes guidées : (1) connexion et configuration,
(2) réponse au questionnaire, (3) obtention du score, (4) analyse des écarts et
feuille de route, (5) vérification de la conformité et export PDF.

\section{Tableau de bord des résultats}
L'utilisateur visualise l'indice global de maturité (en \% et sur 5), un radar
comparant le score courant à la cible, le détail par dimension, ainsi que les
indicateurs clés (indicateurs notés, ignorés, conformité BCT, écarts critiques).

\section{Administration du contenu}
L'administrateur dispose d'un back-office lui permettant de faire évoluer
intégralement le référentiel sans intervention de développement, garantissant la
pérennité de l'outil.

\section{Cohérence et robustesse}
Les libellés et les figures sont dérivés du contenu vivant : le nombre
d'indicateurs et de dimensions, le radar et les barres s'adaptent automatiquement
aux modifications apportées par l'administrateur.

% ===================================================================
\chapter{Conclusion et perspectives}

\section{Bilan}
DataPilot répond à un besoin concret du secteur bancaire tunisien : disposer d'un
instrument fiable, reproductible et conforme pour piloter la maturité de la
gestion des données. La solution livrée couvre l'évaluation, le scoring, l'analyse
d'écarts, la conformité et l'export de rapports, le tout sécurisé par une
authentification sur invitation et une gestion fine des droits.

\section{Perspectives}
\begin{itemize}[noitemsep]
  \item \textbf{Centralisation des soumissions} : stockage serveur des
  évaluations pour une vue consolidée et des comparaisons inter-banques ;
  \item \textbf{Bibliothèque de recommandations} éditable par l'administrateur ;
  \item \textbf{Tableau de bord analytique} pour le suivi dans le temps ;
  \item \textbf{Déploiement on-premise} pour les données bancaires sensibles ;
  \item \textbf{Références réglementaires précises} (article par article) dans le
  rapport de conformité.
\end{itemize}

% ===================================================================
\begin{thebibliography}{9}
\addcontentsline{toc}{chapter}{Bibliographie}
\bibitem{bct} Banque Centrale de Tunisie, \emph{Circulaire N\textsuperscript{o}2025-08
relative à la gouvernance et à la gestion des données}, 2025.
\bibitem{bcbs} Comité de Bâle sur le contrôle bancaire, \emph{BCBS 239 —
Principles for effective risk data aggregation and risk reporting}, 2013.
\bibitem{cmmi} CMMI Institute, \emph{Capability Maturity Model Integration},
ISACA.
\bibitem{react} Meta, \emph{React — A JavaScript library for building user
interfaces}, \url{https://react.dev}.
\bibitem{supabase} Supabase, \emph{The open source Firebase alternative},
\url{https://supabase.com}.
\end{thebibliography}

\end{document}
